\chapter{Минимальное действие петли Тёплица и вариационный разрыв}

Ретроспективный аудит строгих блоков проекта выделяет общий принцип:
действие допустимо, когда оно строится одним оператором, одним
невзвешенным следом и положительным квадратом. Так устроен успешный
связующий функционал $\tau_{300}(([d,d^*]-h)^2)$; обычное спектральное
действие и произвольные центрированные потенциалы ранее не прошли тест
единственности весов.

Для явного унитария $V_{15}$ минимальный аналог этого принципа есть
\begin{equation}
 S_{\rm loop}(V)=\tau_{210}\!\left([N,V]^*[N,V]\right).
 \label{eq:v5-loop-action}
\end{equation}
Он использует только уже выведенные оператор числа, обменный унитарий и
нормированный след. Это стандартная энергия Дирихле петли в унитарной
группе; её критические точки являются гармоническими петлями.

Если собственные winding-числа одной ветви равны $k_1,\ldots,k_{105}$,
то
\begin{equation}
 S_+=\sum_{a=1}^{105}k_a^2,
 \qquad \sum_a k_a=15.
\end{equation}
Целочисленная выпуклость даёт единственный тип минимума:
пятнадцать чисел равны единице, остальные девяносто равны нулю. Поэтому
$zq_0+1-q_0$ является глобальным минимумом в winding-секторе пятнадцать
с точностью до унитарного сопряжения. Сопряжённая ветвь даёт тот же вклад.
После нормировки
\begin{equation}
 S_{\rm loop}(V_{15})=\frac{30}{210}=\frac17.
\end{equation}
Вторая вариация неотрицательна; нули являются касательными к орбите
сопряжения и общим фазовым симметриям.

Этот результат согласуется с онтологией первичности перехода: действие
оценивает стоимость изменения локального шага, а частица пока читается
как минимальный дефект композиции переходов.

\section{Почему это ещё не физическое действие частицы}

Петля $V_{15}(z)$ зависит от внутренней окружности, но не от
$x\in\mathbb R^3$. Поэтому \eqref{eq:v5-loop-action} не создаёт
пространственный профиль и не доказывает локализацию.

Более того, при $N_R=N/R$ получаем либо $S\sim R^{-2}$ при фиксированной
мере, либо $S\sim R^{-1}$ после включения длины окружности. Положительной
стационарной точки по $R$ нет. Это повторяет известный масштабный предел:
квадратичный sigma-model член в трёх измерениях сам не стабилизирует
солитон; в моделях Скирма--Фаддеева баланс требует члена иного порядка по
производным или потенциала. Литература даёт такие механизмы, но проект
пока не выводит их коэффициенты из $N$, $V_{15}$ и общего следа.

\begin{theorem}[Частичное вариационное замыкание]
Функционал $S_{\rm loop}$ канонически выбирает $V_{15}$ как минимальную
гармоническую петлю и воспроизводит вес $1/7$. Он стабилизирует внутренний
winding-класс, но не конечный пространственный размер, массу или энергию
частицы.
\end{theorem}

\begin{protocol}
\begin{enumerate}
 \item один оператор и один след: \textbf{да};
 \item внешний потенциал: \textbf{не введён};
 \item минимум фиксированного winding: \textbf{$V_{15}$};
 \item нормированное действие: \textbf{$1/7$};
 \item положительность второй вариации по петле: \textbf{да, modulo orbit};
 \item стабилизация радиуса: \textbf{нет};
 \item трёхмерная локализация: \textbf{нет};
 \item физическая масса: \textbf{не получена};
 \item следующий гейт: \textbf{пространственное продолжение и баланс Деррика}.
\end{enumerate}
\end{protocol}

Машинный сертификат сохранён в
\path{s2t_v5_toeplitz_parent_action_variational_gap_gate_results.json}.